\subsubsection{Crystals}
\label{sec:crystals}

A crystal tokenizer maps an infinite periodic arrangement to finite, addressable units. An ordered structure can be written as $x=(L,A,F)$, with lattice $L$, element identities $A=(a_i)_{i=1}^{n}$, and fractional coordinates $F=(f_i)_{i=1}^{n}$; its sites are $L(f_i+k)$ for $k\in\mathbb{Z}^{3}$. Atom permutations, origin shifts, and equivalent cell choices change this record while preserving the crystal. Following the survey's three-module definition, we compare the represented object, token units and vocabulary, construction of tokenizer parameters, and the algorithms that encode and decode a structure. Table~\ref{tab:crystal-tokenization} applies these questions to representative methods.

\textbf{Representation.}
Crystal representations retain different parts of composition, geometry, connectivity, and symmetry. CGCNN represents atoms and periodic neighbors as a multigraph for property prediction~\citep{xie2018cgcnn}; coGN supports symmetry-reduced graphs over asymmetric-unit sites and their periodic images~\citep{ruff2024cogn}. DiffCSP retains the lattice and fractional coordinates as explicit geometric variables~\citep{jiao2023diffcsp}. CrystaLLM and Mat2Seq serialize geometric structure records~\citep{antunes2024crystallm,yan2024mat2seq}, whereas SLICES stores composition and periodic connectivity, and SLICES-PLUS appends spatial symmetry operations~\citep{xiao2023slices,wang2025slicesplus}.

Symmetry reduction replaces repeated atoms with orbit representatives. CrystalFormer, PLaID++, and Matra-Genoa retain representative-site coordinates and lattice parameters, allowing orbit expansion into a complete cell~\citep{cao2025crystalformer,xu2026plaidpp,debreuck2026matragenoa}. Coordinate-free Wyckoff descriptions instead group geometric realizations into one decorated symmetry prototype. Goodall et al. use this coarse graining for stability prediction~\citep{goodall2022coarse}; WyFormer uses it for generation~\citep{kazeev2025wyformer}. CDVAE learns a continuous latent representation from which a neural decoder reconstructs the structure~\citep{xie2022cdvae}. These choices determine whether decoding recovers an explicit geometry, infers one from connectivity, or supplies a member of a prototype family.

\textbf{Token units and vocabulary.}
Crystal sites, symmetry orbits, and graph edges become addressable units through vocabulary lookup or continuous feature maps. CrystaLLM defines a symbolic vocabulary of CIF tags, element and space-group symbols, digits, and punctuation. Mat2Seq uses keywords and numerical symbols over a reduced-cell description. PLaID++ passes Wyckoff-field text to a pretrained language tokenizer, so one crystallographic site can occupy several text tokens. SLICES uses element symbols, node indices, and integer cell-translation labels; SLICES-PLUS adds symbols for rotation matrices and fractional translations. Their vocabulary entries describe syntax, chemistry, and relations at different scales.

WyFormer binds an element, site-symmetry category, and enumeration index into a site token, accompanied by a global space-group token. The enumeration distinguishes Wyckoff positions with the same site symmetry within a space group. CrystalFormer uses discrete element and Wyckoff categories together with continuous positional and lattice variables. Matra-Genoa likewise combines categorical fields with continuous geometry, using Gaussian embeddings for continuous values. Thus, a finite vocabulary may cover only part of a crystal record. Atom and edge units in graph methods carry continuous geometric features, while CDVAE uses a continuous latent space rather than a finite codebook.

\textbf{Tokenizer training algorithm.}
Predefined tokenizers construct $\theta_0$ from crystallographic tables, neighborhood rules, serialization grammar, and numerical precision. CGCNN graph construction, Mat2Seq cell reduction, SLICES graph serialization, and Wyckoff-site identification follow such rules. They require no corpus-fitted merge table. CrystalFormer, WyFormer, and Matra-Genoa train embeddings and generative models over prescribed fields; their model training does not learn the field boundaries. PLaID++ inherits its language vocabulary and segmentation rules from the base model, then fine-tunes the model on serialized crystals.

CDVAE learns its encoding and reconstruction maps jointly. A periodic graph encoder predicts a latent Gaussian distribution; training combines a KL penalty with objectives for atom count, composition, lattice parameters, and denoising atom types and positions. Reconstruction losses shape the information stored in the latent. SLICES occupies another point in the design space: encoding uses fixed graph rules, while its reconstruction pipeline reuses a pretrained interatomic potential after geometric optimization. Tokenizer training can therefore consist of fixed parameter construction, reuse of an existing tokenizer or decoder, or joint fitting of an encoder--decoder pair.

\textbf{Encoding algorithm.}
For graph inputs, encoding selects periodic neighbors and constructs atom and edge attributes. CGCNN expands neighbor distances into Gaussian features in its released implementation; coGN supports reduction to asymmetric-unit nodes, retains image transformations and multiplicities, and builds nested line-graph relations for local geometry. DiffCSP extracts $(L,A,F)$ directly, with fractional coordinates treated periodically. CDVAE then adds a learned graph-to-latent mapping when a compressed latent is the output of encoding.

Sequence encoders additionally choose a cell, an ordering, and a numerical representation. CrystaLLM standardizes a CIF and applies lexical token lookup. Mat2Seq performs Niggli reduction, selects an origin, orders irreducible sites, and serializes their geometry and symmetry information; its construction targets rigid and periodic invariance. SLICES identifies neighbors with EconNN and writes a labeled quotient graph. SLICES-PLUS adds a symbolic sequence of symmetry operations. Wyckoff encoders identify orbits and attach elements: CrystalFormer orders Wyckoff letters and representative coordinates, WyFormer maps positions to site-symmetry and enumeration categories, and PLaID++ refines a cell before writing its Wyckoff text. Matra-Genoa serializes categorical and continuous fields into an invertible mixed sequence. Ordering and symmetry tolerance enter at this stage, before a generative model predicts any tokens.

\textbf{Decoding algorithm.}
Explicit geometric records support parsing and coordinate recovery. CrystaLLM concatenates lexical tokens into CIF text for parsing. Mat2Seq recovers lattice parameters and expands irreducible sites using symmetry operations. CrystalFormer and Matra-Genoa recover representative geometry and apply the corresponding Wyckoff operations; PLaID++ converts language tokens to text, parses its fields, and expands sites through PyXtal. Stored precision determines geometric resolution: Mat2Seq rounds real values to four decimal places, while PLaID++ uses two for lattice parameters and three for fractional coordinates.

Topology and latent representations require additional geometry construction. SLICES' SLI2Cry routine embeds the periodic graph, optimizes chemically informed geometry, and refines it with an interatomic potential; SLICES-PLUS augments the graph representation with symmetry information. WyFormer decodes to a protostructure, after which PyXtal and geometry optimization or a separate geometric model supply coordinates and lattice values. CDVAE predicts global structure attributes from its latent and uses Langevin denoising to reconstruct atoms and positions. CGCNN and coGN specify property-oriented graph encodings without a structure-reconstruction decoder; DiffCSP's state already contains explicit geometry, and its reverse diffusion learns a distribution over that state. Figure~\ref{fig:crystal-tokenization-paths} contrasts these encoding and decoding paths.

\textbf{Coverage and evaluation.}
Defects change the object that encoding must preserve. A periodic supercell fixes both defect concentration and repeat distance; symmetry lowering increases the number of independent sites and reduces orbit-based compression. Partial occupancy describes an average disordered structure, while CrystaLLM's training data contain ordered structures~\citep{antunes2024crystallm}. Element coverage, occupancy conventions, cell size, and symmetry tolerance therefore enter the input definition.

Tokenizer evaluation can follow the same module boundaries. For \texttt{train}, report fitted or reused components and the data used to obtain them. For \texttt{encode}, measure token and relation counts, numerical precision, cost, and consistency under equivalent cells, origin shifts, and atom permutations. For \texttt{decode}, measure structural matching and geometric error at a stated tolerance, separating record recovery from geometry inferred by optimization or a learned decoder. Changes in composition, defects, or free coordinates should remain distinguishable when they belong to the target. Diffraction supplies a linked measurement for evaluating reconstructed crystals under specified experimental conditions.

\begin{table}[H]
\centering
\caption{Features of representative crystal tokenizations. Checkmarks indicate neural latent encoding (N), explicit complete geometry (G), and explicit crystallographic symmetry information (S). Type: D, discrete; M, mixed discrete--continuous; C, continuous latent.}
\label{tab:crystal-tokenization}
\small
\setlength{\tabcolsep}{4pt}
\renewcommand{\arraystretch}{1.17}
\begin{tabularx}{\textwidth}{@{}L{0.16\textwidth}YccccL{0.17\textwidth}L{0.18\textwidth}@{}}
\toprule
Method & Units & Type & N & G & S & Tokenizer training & Reconstruction \\
\midrule
CGCNN~\citep{xie2018cgcnn} & Atoms, periodic edges & M & -- & -- & -- & Fixed graph rules & Not specified \\
coGN~\citep{ruff2024cogn} & Reduced sites, image edges & M & -- & -- & $\checkmark$ & Fixed graph rules & Not specified \\
DiffCSP~\citep{jiao2023diffcsp} & Elements, coordinates, lattice & M & -- & $\checkmark$ & -- & Fixed fields & Explicit state \\
CDVAE~\citep{xie2022cdvae} & Crystal latent & C & $\checkmark$ & -- & -- & Joint encoder--decoder & Learned, approximate \\
CrystaLLM~\citep{antunes2024crystallm} & CIF symbols and digits & D & -- & $\checkmark$ & $\checkmark$ & Fixed vocabulary & CIF parsing \\
Mat2Seq~\citep{yan2024mat2seq} & Reduced-cell and site fields & D & -- & $\checkmark$ & $\checkmark$ & Fixed rules & Symmetry expansion \\
SLICES~\citep{xiao2023slices} & Elements, indices, cell shifts & D & -- & -- & -- & Fixed encoder; reused potential & Graph embedding and optimization \\
SLICES-PLUS~\citep{wang2025slicesplus} & Graph and symmetry symbols & D & -- & -- & $\checkmark$ & Fixed encoder & Graph reconstruction \\
CrystalFormer~\citep{cao2025crystalformer} & Wyckoff sites and geometry & M & -- & $\checkmark$ & $\checkmark$ & Fixed fields & Symmetry expansion \\
WyFormer~\citep{kazeev2025wyformer} & Decorated Wyckoff tuples & D & -- & -- & $\checkmark$ & Fixed tuples & Prototype; geometry added later \\
PLaID++~\citep{xu2026plaidpp} & Text over Wyckoff fields & D & -- & $\checkmark$ & $\checkmark$ & Reused language tokenizer & Parsing and symmetry expansion \\
Matra-Genoa~\citep{debreuck2026matragenoa} & Discrete fields and geometry & M & -- & $\checkmark$ & $\checkmark$ & Fixed fields & Mixed-field decoding \\
\bottomrule
\end{tabularx}
\vspace{0.4em}
\begin{minipage}{\textwidth}\footnotesize
N refers to the structural latent encoder, rather than embeddings in the consuming model. G requires both lattice and sufficient positional variables; S can be present without G. The DiffCSP row describes its explicit state representation, and the CDVAE row its encoded latent. A dash marks absence in the selected representation.
\end{minipage}
\end{table}

\begin{figure}[H]
\centering
\resizebox{\textwidth}{!}{%
\begin{tikzpicture}[y=0.73cm,font=\footnotesize,>=Latex,
  arrow/.style={->,line width=0.9pt,draw=darkgray!80},
  chip/.style={draw=darkgray!30,rounded corners=2pt,fill=white,inner xsep=5pt,inner ysep=2pt,font=\footnotesize},
  note/.style={align=center,font=\footnotesize,text=darkgray},
  cs/.style={circle,draw=deepblue!90!black,fill=midblue,minimum size=7pt,inner sep=0pt},
  cl/.style={circle,draw=warm!90!black,fill=warm!75,minimum size=9pt,inner sep=0pt}]
% Shared illustrative input: one CsCl-type cell.
\begin{scope}[shift={(4.05,13.05)},scale=0.90]
\draw[darkgray!50,dashed] (0,0)--(0.55,0.45)--(1.9,0.45) (0.55,0.45)--(0.55,1.8);
\draw[darkgray!65] (0,0) rectangle (1.35,1.35) (0,1.35)--(0.55,1.8)--(1.9,1.8)--(1.35,1.35) (1.35,0)--(1.9,0.45)--(1.9,1.8);
\foreach \x/\y in {0/0,1.35/0,0/1.35,1.35/1.35,0.55/0.45,1.9/0.45,0.55/1.8,1.9/1.8}{\node[cs] at (\x,\y) {};}
\node[cl] at (0.95,0.9) {};
\end{scope}
\node[anchor=west,font=\bfseries\normalsize] at (6.2,14.15) {One periodic crystal};
\node[anchor=west,text=darkgray] at (6.2,13.65) {CsCl-type cell: $x=(L,A,F)$};
\node[cs] at (6.35,13.1) {}; \node[anchor=west,font=\footnotesize] at (6.55,13.1) {Cs};
\node[cl] at (7.6,13.1) {}; \node[anchor=west,font=\footnotesize] at (7.8,13.1) {Cl};
\draw[arrow] (7,12.8)--(7,12.25);
\node[note] at (7,11.96) {Four choices of what becomes a token};

% Four illustrated panels.
\foreach \x/\y/\c in {0/5.65/deepblue,7.2/5.65/warm,0/-0.65/deepblue,7.2/-0.65/warm}{
 \fill[\c!3] (\x,\y) rectangle +(6.8,6.0);
 \draw[\c!35,line width=0.6pt] (\x,\y) rectangle +(6.8,6.0);
}

% A: full geometric record -> lexical pieces.
\node[anchor=west,font=\bfseries] at (0.3,11.22) {A\quad Explicit geometry};
\node[anchor=west,font=\footnotesize,text=darkgray] at (0.3,10.82) {CrystaLLM; geometric text in Mat2Seq};
\draw[fill=white,draw=deepblue!40] (0.65,8.6) rectangle (3.75,10.4);
\foreach \y in {10.05,9.7,9.35,9.0}{\draw[deepblue!15] (0.85,\y)--(3.55,\y);}
\node[anchor=west,font=\ttfamily\footnotesize] at (0.92,10.06) {a = b = c};
\node[anchor=west,font=\ttfamily\footnotesize] at (0.92,9.7) {Cs  0  0  0};
\node[anchor=west,font=\ttfamily\footnotesize] at (0.92,9.34) {Cl .5 .5 .5};
\node[anchor=west,font=\footnotesize,text=darkgray] at (4.02,9.5) {serialize};
\draw[arrow] (2.2,8.48)--(2.2,8.02);
\node[chip,fill=deepblue!10] at (0.9,7.65) {\texttt{Cs}};
\node[chip] at (1.7,7.65) {\texttt{0}};
\node[chip] at (2.4,7.65) {\texttt{0}};
\node[chip] at (3.1,7.65) {\texttt{0}};
\node[chip,fill=warm!12] at (4.0,7.65) {\texttt{Cl}};
\node[chip] at (4.9,7.65) {\texttt{.5}};
\node at (5.8,7.65) {$\cdots$};
\node[note,text width=6.1cm] at (3.4,6.8) {Token units: symbols / numbers\\\textbf{Decode:} parse the stored geometry};
\node[note,font=\scriptsize] at (3.4,6.03) {Training: fixed serialization / vocabulary};

% B: many atoms -> two orbit representatives.
\node[anchor=west,font=\bfseries] at (7.5,11.22) {B\quad Symmetry orbits};
\node[anchor=west,align=left,font=\footnotesize,text=darkgray] at (7.5,10.55) {WyFormer; with geometry:\\PLaID++, CrystalFormer};
\begin{scope}[shift={(7.85,8.76)}]
\foreach \x/\y in {0/0,1.0/0,0/1.0,1.0/1.0}{\node[cs] at (\x,\y) {};}
\node[cl] at (0.5,0.5) {};
\draw[darkgray!30] (0,0) rectangle (1,1);
\draw[arrow] (1.35,.5)--(2.05,.5);
\node[cs,minimum size=13pt] at (2.5,.85) {};
\node[cl,minimum size=13pt] at (2.5,.13) {};
\node[anchor=west,font=\footnotesize] at (2.8,.85) {Cs orbit: $1a$};
\node[anchor=west,font=\footnotesize] at (2.8,.13) {Cl orbit: $1b$};
\end{scope}
\node[chip,fill=warm!12] at (8.8,7.95) {SG 221};
\node[chip,fill=deepblue!10] at (10.55,7.95) {Cs + site ID};
\node[chip,fill=warm!12] at (12.55,7.95) {Cl + site ID};
\node[note,text width=6.1cm] at (10.6,6.94) {\textbf{With geometry:} expand the orbits\\\textbf{Without geometry:} add it later};
\node[note,font=\scriptsize] at (10.6,6.03) {Training: fixed orbit rules / reused text tokenizer};

% C: periodic neighbors and cell shifts -> edge records.
\node[anchor=west,font=\bfseries] at (0.3,4.92) {C\quad Periodic connectivity};
\node[anchor=west,font=\footnotesize,text=darkgray] at (0.3,4.52) {SLICES / SLICES-PLUS};
\draw[darkgray!35,dashed] (0.65,2.35) rectangle (2.7,4.1);
\draw[darkgray!20,dashed] (2.7,2.35) rectangle (4.75,4.1);
\node[cs] (g1) at (1.05,3.7) {};\node[cl] (g2) at (2.1,2.85) {};
\node[cs] (g3) at (3.1,3.7) {};\node[cl] (g4) at (4.15,2.85) {};
\draw[deepblue!70,line width=1.1pt] (g1)--(g2) (g3)--(g4);
\draw[warm,line width=1.2pt] (g2)--(g3);
\node[font=\footnotesize,fill=white,inner sep=2pt] at (3.02,3.14) {$+a$};
\node[note] at (5.5,3.15) {cell shift\\is retained};
\node[chip] at (3.4,1.68) {Cs, Cl;\quad $(u,v,n_x,n_y,n_z)$};
\node[note,text width=6cm] at (3.4,.83) {Tokens store periodic connections\\\textbf{Decode:} infer geometry from the graph};
\node[note,font=\scriptsize] at (3.4,-.23) {Training: fixed graph rules + reused potential};

% D: learned compression into real-valued latent coordinates.
\node[anchor=west,font=\bfseries] at (7.5,4.92) {D\quad Learned latent};
\node[anchor=west,font=\footnotesize,text=darkgray] at (7.5,4.52) {CDVAE};
\node[cs] (l1) at (7.95,3.7) {};\node[cl] (l2) at (8.8,3.5) {};
\node[cs] (l3) at (8.3,2.7) {};\node[cl] (l4) at (9.05,2.45) {};
\draw[darkgray!45] (l1)--(l2)--(l4)--(l3)--(l1) (l2)--(l3);
\draw[fill=deepblue!10,draw=deepblue!55] (9.45,3.9)--(9.95,3.9)--(10.35,3.25)--(9.95,2.6)--(9.45,2.6)--cycle;
\node[font=\footnotesize,rotate=90] at (9.8,3.25) {encoder};
\draw[arrow] (9.02,3.25)--(9.4,3.25);
\foreach \x/\h/\c in {10.75/.7/20,11.05/1.1/65,11.35/.4/35,11.65/.9/80}{\fill[midblue!\c] (\x,2.75) rectangle +(0.22,\h);}
\node[note] at (11.3,2.42) {$z\in\mathbb{R}^{d}$};
\draw[arrow] (12,3.25)--(12.6,3.25);
\node[cl] at (13.05,3.65) {};\node[cs] at (13.5,2.9) {};\node[cs] at (12.9,2.6) {};
\draw[darkgray!35] (13.05,3.65)--(13.5,2.9)--(12.9,2.6)--cycle;
\node[note] at (13.1,2.08) {denoise};
\node[note,text width=6cm] at (10.6,.83) {A learned vector replaces the input\\\textbf{Decode:} neural reconstruction of a crystal};
\node[note,font=\scriptsize] at (10.6,-.23) {Training: jointly learn encoder and decoder};
\end{tikzpicture}%
}
\caption{Four representative tokenization families illustrated from a CsCl-type periodic structure. Colored circles distinguish elements; graph sketches are schematic, and chips summarize fields rather than literal model token streams. Full geometry, orbit decorations, periodic connectivity, and learned latents retain different information and require different decoding algorithms.}
\label{fig:crystal-tokenization-paths}
\end{figure}
\clearpage
